\documentclass[12pt, a4paper]{article}

% ── Shared preamble (src/lecturenotes.sty) ────────────────────────────────────
%   Packages, colors, boxes, \sol/\soleq and math macros all live in
%   src/*.sty. \input@path points LaTeX two folders up to find them.
\makeatletter
\def\input@path{{../../}}
\makeatother
\usepackage{lecturenotes}

% ── Document ──────────────────────────────────────────────────────────────────
\begin{document}

\begin{center}
  {\LARGE \textbf{Section 3.4: Backpropagation}} \\[0.3em]
  {\large\itshape Computing Gradients with the Chain Rule} \\[0.5em]
  {\large AI-IM Class --- Lecture Notes}%
  \ifstudentmode
    \quad {\normalsize\color{gray}[Student Copy]}%
  \fi
\end{center}

\vspace{1em}
\hrule
\vspace{1.5em}

% ══════════════════════════════════════════════════════════════════════════════
\section{The Goal}

From Section 3.2 (Gradient Descent), we know that to update a weight (and bias) $w$
we need its gradient:
\[
  w \;\longleftarrow\; w - \eta \,\pd{J}{w}.
\]

From Section 3.3 (Forward Pass) we computed a prediction $\yhat = 0.637$
and a loss $J = 0.132$ for the input $(x_1, x_2) = (0, 1)$.

\textbf{Now we need $\partial J / \partial w$ for every weight in the
  network.}
The difficulty is that the loss depends on a weight only \emph{indirectly}:
$w \to z \to h \to \cdots \to \yhat \to J$.
Each arrow is a function, and we need to differentiate through all of them.
The tool that makes this systematic is the \textbf{chain rule}.

% ══════════════════════════════════════════════════════════════════════════════
\section{The Chain Rule}

If $y$ depends on $u$, and $u$ depends on $x$, then:
\[
  \frac{dy}{dx} = \frac{dy}{du} \cdot \frac{du}{dx}.
\]

This extends to any chain of functions.
For three nested functions $y = f(g(h(x)))$, let $v = h(x)$ and $u = g(v)$:

\[
  \frac{dy}{dx}
  = \frac{dy}{du} \cdot \frac{du}{dv} \cdot \frac{dv}{dx}.
\]

\begin{example}
  Let $y = (2x+1)^2$.  Write $u = 2x+1$ so that $y = u^2$:
  \[
    \frac{dy}{dx}
    = \underbrace{\frac{dy}{du}}_{2u}
    \cdot
    \underbrace{\frac{du}{dx}}_{2}
    = 2(2x+1) \cdot 2
    = \sol[4cm]{4(2x+1)}.
  \]
\end{example}

\begin{notebox}
  Backpropagation is simply the chain rule applied \emph{backwards} through
  the network --- one layer at a time.  We start from $\partial J/\partial
    \yhat$ (easy to compute) and work leftward, accumulating the product at
  each layer.
\end{notebox}

% ══════════════════════════════════════════════════════════════════════════════
\section{What We Know from Section 3.3}
The network has the following (current) parameters:
\[
  \bW^{(1)} = \begin{bmatrix} 0.5 & -0.5 \\ 0.5 & \phantom{-}0.5 \end{bmatrix},
  \quad
  \bb^{(1)} = \begin{bmatrix} 0.5 \\ 0.0 \end{bmatrix},
  \quad
  \bW^{(2)} = \begin{bmatrix} 0.5 & 0.5 \end{bmatrix},
  \quad
  b^{(2)} = 0.0.
\]

\begin{figure}[H]
  \centering
  \begin{tikzpicture}[
      neuron/.style={circle, draw=mainblue, fill=boxbg,
          minimum size=1.1cm, thick, font=\small},
      preact/.style={circle, draw=mainblue, fill=white,
          minimum size=1.1cm, thick, font=\small},
      input/.style={circle, draw=gray!60, fill=gray!15,
          minimum size=1.1cm, thick, font=\small},
      unitbox/.style={draw=gray!60, dashed, rounded corners=6pt},
      >=Stealth, node distance=0pt
    ]
    % ── Nodes ──────────────────────────────────────────────────────────────────
    \node[input]   (x1) at (0,    1.4) {$x_1$};
    \node[input]   (x2) at (0,   -1.4) {$x_2$};

    \node[preact]  (z1) at (3.6,  1.4) {$z^{(1)}_1$};
    \node[neuron]  (h1) at (5.6,  1.4) {$h_1$};
    \node[preact]  (z2) at (3.6, -1.4) {$z^{(1)}_2$};
    \node[neuron]  (h2) at (5.6, -1.4) {$h_2$};

    \node[preact]  (zo) at (9.2,  0)   {$z^{(2)}$};
    \node[neuron]  (yh) at (11.2, 0)   {$\yhat$};

    % ── Dashed boxes: one box = one neuron ─────────────────────────────────────
    \draw[unitbox] (2.85,  0.65) rectangle (6.35,  2.15);
    \draw[unitbox] (2.85, -2.15) rectangle (6.35, -0.65);
    \draw[unitbox] (8.45, -0.75) rectangle (11.95, 0.75);

    % ── Connections: input → hidden pre-activations ────────────────────────────
    \draw[->, thick] (x1) -- node[above, font=\scriptsize, pos=0.5]
    {$w^{(1)}_{11}$} (z1);
    \draw[->, thick] (x1) -- node[sloped, above, font=\scriptsize, pos=0.28]
    {$w^{(1)}_{21}$} (z2);
    \draw[->, thick] (x2) -- node[sloped, below, font=\scriptsize, pos=0.28]
    {$w^{(1)}_{12}$} (z1);
    \draw[->, thick] (x2) -- node[below, font=\scriptsize, pos=0.5]
    {$w^{(1)}_{22}$} (z2);

    % ── Activation arrows z → h ────────────────────────────────────────────────
    \draw[->, thick, mainblue] (z1) -- node[above, font=\small] {$\sigmoid$} (h1);
    \draw[->, thick, mainblue] (z2) -- node[above, font=\small] {$\sigmoid$} (h2);
    \draw[->, thick, mainblue] (zo) -- node[above, font=\small] {$\sigmoid$} (yh);

    % ── Connections: hidden → output pre-activation ────────────────────────────
    \draw[->, thick] (h1) -- node[above right, font=\scriptsize, pos=0.45]
    {$w^{(2)}_1$} (zo);
    \draw[->, thick] (h2) -- node[below right, font=\scriptsize, pos=0.45]
    {$w^{(2)}_2$} (zo);

    % ── Bias annotations ───────────────────────────────────────────────────────
    \node[font=\scriptsize, gray!80!black] at (3.6,  2.45) {$+\,b^{(1)}_1$};
    \node[font=\scriptsize, gray!80!black] at (3.6, -2.45) {$+\,b^{(1)}_2$};
    \node[font=\scriptsize, gray!80!black] at (9.2,  1.05) {$+\,b^{(2)}$};

    % ── Layer labels ───────────────────────────────────────────────────────────
    \node[font=\small\itshape, gray] at (0,    3.1) {Input layer};
    \node[font=\small\itshape, gray] at (4.6,  3.1) {Hidden layer};
    \node[font=\small\itshape, gray] at (10.2, 3.1) {Output layer};
    \node[font=\small\itshape, gray] at (13.6, 3.1) {Loss};

    % ── Loss node and target ───────────────────────────────────────────────────
    \node[circle, draw=red!70!black, fill=warnbg, minimum size=1.1cm,
      thick, font=\small] (J) at (13.6, 0) {$J$};
    \node[input] (y) at (13.6, -2.0) {$y$};
    \draw[->, thick, red!70!black] (yh) -- (J);
    \draw[->, thick, red!70!black] (y)  -- (J);
    \node[font=\scriptsize, red!70!black] at (13.6, 0.95) {$J = (y - \yhat)^2$};
    \node[font=\scriptsize, gray!80!black, left=0.1cm of y] {target};
  \end{tikzpicture}
  \caption{The complete forward pass.}
  \label{fig:forward-pass-z-with-loss}
\end{figure}

\begin{center}
  \renewcommand{\arraystretch}{1.4}
  \begin{tabular}{lll}
    \toprule
    Quantity                & Symbol         & Value      \\
    \midrule
    Input                   & $(x_1,\, x_2)$ & $(0,\; 1)$ \\
    Target                  & $y$            & $1$        \\
    \midrule
    Hidden pre-activation 1 & $z^{(1)}_1$    & $0.0$      \\
    Hidden pre-activation 2 & $z^{(1)}_2$    & $0.5$      \\
    Hidden activation 1     & $h_1$          & $0.500$    \\
    Hidden activation 2     & $h_2$          & $0.622$    \\
    \midrule
    Output pre-activation   & $z^{(2)}$      & $0.561$    \\
    Prediction              & $\yhat$        & $0.637$    \\
    \midrule
    Loss                    & $J$            & $0.132$    \\
    \bottomrule
  \end{tabular}
\end{center}


\begin{notebox}
  We also need the sigmoid derivatives $\sigmoid'(z) = \sigmoid(z)(1-\sigmoid(z))$
  at the values we computed.  We can use Python to evaluate them:

\end{notebox}
\begin{codebox}
def sigmoid_p(x): # derivative of sigmoid
    s = sigmoid(x)
    return s * (1 - s)
\end{codebox}


% ══════════════════════════════════════════════════════════════════════════════
\section{Backward Pass: Output Layer}

\begin{figure}[H]
  \centering
  \begin{tikzpicture}[
      neuron/.style={circle, draw=mainblue, fill=boxbg,
          minimum size=1.1cm, thick, font=\small},
      preact/.style={circle, draw=mainblue, fill=white,
          minimum size=1.1cm, thick, font=\small},
      input/.style={circle, draw=gray!60, fill=gray!15,
          minimum size=1.1cm, thick, font=\small},
      param/.style={rectangle, rounded corners=3pt, draw=orange!80!black,
          fill=notebg, thick, minimum size=0.75cm, font=\small},
      unitbox/.style={draw=gray!60, dashed, rounded corners=6pt},
      >=Stealth, node distance=0pt
    ]
    % ── Nodes ──────────────────────────────────────────────────────────────────
    \node[preact]  (zo) at (9.2,  0)   {$z^{(2)}$};
    \node[neuron]  (yh) at (11.2, 0)   {$\yhat$};
    \node[circle, draw=red!70!black, fill=warnbg, minimum size=1.1cm,
      thick, font=\small] (J) at (13.6, 0) {$J$};
    \node[input] (y) at (13.6, -2.0) {$y$};

    % ── Output neuron box ──────────────────────────────────────────────────────
    \draw[unitbox] (8.45, -0.75) rectangle (11.95, 0.75);

    % ── Incoming arrows from the hidden layer (h nodes omitted) ────────────────
    \draw[->, thick] (5.6,  1.4) -- (zo);
    \draw[->, thick] (5.6, -1.4) -- (zo);

    % ── Parameters of the output neuron, linked to z by dotted lines ──────────
    \node[param] (w1) at (6.9,  2.3) {$w^{(2)}_1$};
    \node[param] (w2) at (6.9, -2.3) {$w^{(2)}_2$};
    \node[param] (b2) at (9.2,  2.3) {$b^{(2)}$};
    \draw[dotted, very thick, orange!80!black] (w1) -- (zo);
    \draw[dotted, very thick, orange!80!black] (w2) -- (zo);
    \draw[dotted, very thick, orange!80!black] (b2) -- (zo);

    % ── Forward arrows ─────────────────────────────────────────────────────────
    \draw[->, thick, mainblue] (zo) -- node[above, font=\small] {$\sigmoid$} (yh);
    \draw[->, thick, red!70!black] (yh) -- (J);
    \draw[->, thick, red!70!black] (y)  -- (J);
    \node[font=\scriptsize, red!70!black] at (13.6, 0.95) {$J = (y - \yhat)^2$};
    \node[font=\scriptsize, gray!80!black, left=0.1cm of y] {target};

    % ── Backward (chain-rule) arrows ───────────────────────────────────────────
    \draw[->, thick, red!70!black, dashed] (J.south west) to[bend left=35]
    node[below, font=\scriptsize] {$\pd{J}{\yhat}$} (yh.south east);
    \draw[->, thick, red!70!black, dashed] (yh.south west) to[bend left=35]
    node[below, font=\scriptsize] {$\pd{\yhat}{z^{(2)}}$} (zo.south east);
  \end{tikzpicture}
  \caption{Zooming in on the output neuron.  The dotted lines show that
    $w^{(2)}_1$, $w^{(2)}_2$ and $b^{(2)}$ affect the loss only through
    $z^{(2)}$, so each gradient is the chain
    $J \to \yhat \to z^{(2)} \to$ parameter.}
  \label{fig:output-chain}
\end{figure}


We work right to left through the computation graph:

\begin{align*}
  J \;\xleftarrow{\partial J/\partial \yhat}\;
  \yhat \;\xleftarrow{\partial \yhat/\partial z^{(2)}}\;
  z^{(2)} & \;\xleftarrow{\partial z^{(2)}/\partial w^{(2)}_j}\;
  w^{(2)}_j \qquad (j = 1, 2)                                    \\[0.4em]
  J \;\xleftarrow{\partial J/\partial \yhat}\;
  \yhat \;\xleftarrow{\partial \yhat/\partial z^{(2)}}\;
  z^{(2)} & \;\xleftarrow{\partial z^{(2)}/\partial b^{(2)}}\;
  b^{(2)}
\end{align*}

The output neuron has three parameters: two weights $w^{(2)}_1, w^{(2)}_2$
and one bias $b^{(2)}$.  All three reach $J$ through $z^{(2)}$
(see Figure~\ref{fig:output-chain}), so their chains share the first two
links and differ only in the last.

% ──────────────────────────────────────────────────────────────────────────────
\begin{stepbox}{Step 1 \quad Gradient of the loss with respect to $\yhat$}

  \[
    J = (y - \yhat)^2
    \implies
    \pd{J}{\yhat}
    = 2(\yhat - y)
    = 2(0.637 - 1)
    = \sol[3cm]{-0.726}
  \]

  The negative sign means: increasing $\yhat$ would \emph{decrease} $J$.
  Since $\yhat < y$, the network is predicting too low.
\end{stepbox}

\vspace{0.8em}

% ──────────────────────────────────────────────────────────────────────────────
\begin{stepbox}{Step 2 \quad Gradient through the output sigmoid}

  \[
    \yhat = \sigmoid(z^{(2)})
    \implies
    \pd{\yhat}{z^{(2)}}
    = \sigmoid'(z^{(2)})
    = \yhat(1-\yhat)
    = 0.637 \times 0.363
    \approx \sol[3cm]{0.231}
  \]

  By the chain rule, the gradient arriving at $z^{(2)}$ is:

  \[
    \pd{J}{z^{(2)}}
    = \pd{J}{\yhat} \cdot \pd{\yhat}{z^{(2)}}
    = (-0.726)(0.231)
    \approx \sol[3cm]{-0.168}
  \]

  We will call this quantity $\delta^{(2)} = -0.168$.
  It is the \textbf{error signal} at the output neuron: how much the loss
  changes per unit change in $z^{(2)}$.
\end{stepbox}

\vspace{0.8em}

% ──────────────────────────────────────────────────────────────────────────────
\begin{figure}[H]
  \centering
  \begin{tikzpicture}[
      preact/.style={circle, draw=mainblue, fill=white,
          minimum size=1.1cm, thick, font=\small},
      param/.style={rectangle, rounded corners=3pt, draw=orange!80!black,
          fill=notebg, thick, minimum size=0.75cm, font=\small},
      linklabel/.style={font=\small, orange!60!black},
      >=Stealth
    ]
    % ── Output pre-activation ──────────────────────────────────────────────────
    \node[preact] (zo) at (6, 0) {$z^{(2)}$};

    % ── Incoming activations (no h nodes, just labelled tails) ────────────────
    \draw[->, thick] (1.8,  0.9) node[left, font=\small] {$h_1$} -- (zo);
    \draw[->, thick] (1.8, -0.9) node[left, font=\small] {$h_2$} -- (zo);

    % ── Parameters and their local derivatives ─────────────────────────────────
    \node[param] (w1) at (3.0,  3.0) {$w^{(2)}_1$};
    \node[param] (w2) at (3.0, -3.0) {$w^{(2)}_2$};
    \node[param] (b2) at (6.0,  3.0) {$b^{(2)}$};
    \draw[dotted, very thick, orange!80!black] (w1) --
    node[linklabel, above right, pos=0.45] {$\pd{z^{(2)}}{w^{(2)}_1}$} (zo);
    \draw[dotted, very thick, orange!80!black] (w2) --
    node[linklabel, below right, pos=0.45] {$\pd{z^{(2)}}{w^{(2)}_2}$} (zo);
    \draw[dotted, very thick, orange!80!black] (b2) --
    node[linklabel, right, pos=0.5, xshift=2pt] {$\pd{z^{(2)}}{b^{(2)}}$} (zo);

    % ── Error signal arriving from the right (Steps 1–2 collapsed) ─────────────
    \draw[->, thick, red!70!black, dashed] (11.4, 0) --
    node[above, font=\small] {$\delta^{(2)} = \pd{J}{z^{(2)}}$}
    node[below, font=\scriptsize, gray!80!black, yshift=-2pt]
      {from $\yhat$ and $J$ (Steps 1--2)} (zo);
  \end{tikzpicture}
  \caption{After Steps 1--2, everything to the right of $z^{(2)}$ collapses
    into the error signal $\delta^{(2)}$.  Each parameter's gradient is
    $\delta^{(2)}$ times the label on its link.}
  \label{fig:output-local}
\end{figure}

\begin{stepbox}{Step 3 \quad Gradients for the output weights}

  Because $z^{(2)} = w^{(2)}_1 h_1 + w^{(2)}_2 h_2 + b^{(2)}$, we have:
  \[
    \pd{z^{(2)}}{w^{(2)}_1} = h_1, \qquad
    \pd{z^{(2)}}{w^{(2)}_2} = h_2, \qquad
    \pd{z^{(2)}}{b^{(2)}} = 1.
  \]

  Applying the chain rule one more time:

  \begin{align*}
    \pd{J}{w^{(2)}_1}
     & = \delta^{(2)} \cdot h_1
    = (-0.168)(0.500)
    = \sol[3cm]{-0.084}         \\[6pt]
    \pd{J}{w^{(2)}_2}
     & = \delta^{(2)} \cdot h_2
    = (-0.168)(0.622)
    \approx \sol[3cm]{-0.105}   \\[6pt]
    \pd{J}{b^{(2)}}
     & = \delta^{(2)} \cdot 1
    = \sol[3cm]{-0.168}
  \end{align*}
\end{stepbox}

\vspace{0.5em}

\begin{notebox}
  \textbf{Pattern observed.}
  For any output weight $w^{(2)}_i$, the gradient is:
  \[
    \pd{J}{w^{(2)}_i} = \delta^{(2)} \cdot h_i
    \qquad \text{(error signal} \times \text{incoming activation)}.
  \]
  This pattern repeats at every layer.
\end{notebox}

% ══════════════════════════════════════════════════════════════════════════════
\section{Backward Pass: Hidden Layer}

\begin{figure}[H]
  \centering
  \begin{tikzpicture}[
      neuron/.style={circle, draw=mainblue, fill=boxbg,
          minimum size=1.1cm, thick, font=\small},
      preact/.style={circle, draw=mainblue, fill=white,
          minimum size=1.1cm, thick, font=\small},
      linklabel/.style={font=\small, orange!60!black},
      deltalabel/.style={font=\small, red!70!black},
      >=Stealth
    ]
    % ── Nodes ──────────────────────────────────────────────────────────────────
    \node[preact] (z1) at (2.4,  1.7) {$z^{(1)}_1$};
    \node[neuron] (h1) at (5.6,  1.7) {$h_1$};
    \node[preact] (z2) at (2.4, -1.7) {$z^{(1)}_2$};
    \node[neuron] (h2) at (5.6, -1.7) {$h_2$};
    \node[preact] (zo) at (10.0, 0)   {$z^{(2)}$};

    % ── Input stubs into the hidden pre-activations (no x nodes) ───────────────
    \draw[->, thick] (0.6,  2.3) -- (z1);
    \draw[->, thick] (0.6,  1.1) -- (z1);
    \draw[->, thick] (0.6, -1.1) -- (z2);
    \draw[->, thick] (0.6, -2.3) -- (z2);

    % ── Activation arrows z → h: σ outside, local derivative inside ────────────
    \draw[->, thick, mainblue] (z1) --
    node[above, font=\small] {$\sigmoid$}
    node[linklabel, below] {$\pd{h_1}{z^{(1)}_1}$} (h1);
    \draw[->, thick, mainblue] (z2) --
    node[below, font=\small] {$\sigmoid$}
    node[linklabel, above] {$\pd{h_2}{z^{(1)}_2}$} (h2);

    % ── Hidden → output: local derivative outside, weight inside ───────────────
    \draw[->, thick] (h1) --
    node[linklabel, above right, pos=0.45] {$\pd{z^{(2)}}{h_1}$}
    node[below left, font=\small, pos=0.5] {$w^{(2)}_1$} (zo);
    \draw[->, thick] (h2) --
    node[linklabel, below right, pos=0.45] {$\pd{z^{(2)}}{h_2}$}
    node[above left, font=\small, pos=0.5] {$w^{(2)}_2$} (zo);

    % ── Error signal arriving from the right (Steps 1–2) ───────────────────────
    \draw[->, thick, red!70!black, dashed] (13.4, 0) --
    node[above, font=\small] {$\delta^{(2)} = \pd{J}{z^{(2)}}$}
    node[below, font=\scriptsize, gray!80!black, yshift=-2pt]
      {from $\yhat$ and $J$} (zo);

    % ── Destinations: hidden error signals ─────────────────────────────────────
    \node[deltalabel, above=0.15cm of z1]
    {$\delta^{(1)}_1 = \pd{J}{z^{(1)}_1}$};
    \node[deltalabel, below=0.15cm of z2]
    {$\delta^{(1)}_2 = \pd{J}{z^{(1)}_2}$};
  \end{tikzpicture}
  \caption{The error signal $\delta^{(2)}$ flows backwards through each
    hidden unit.  Multiplying the labels along a path gives
    $\delta^{(1)}_i = \delta^{(2)} \cdot \frac{\partial z^{(2)}}{\partial h_i}
      \cdot \frac{\partial h_i}{\partial z^{(1)}_i}$.}
  \label{fig:hidden-chain}
\end{figure}

To update $\bW^{(1)}$ and $\bb^{(1)}$ we need
$\partial J / \partial w^{(1)}_{ij}$ and $\partial J / \partial b^{(1)}_i$.
Working right to left, the chains are:

\begin{align*}
  J \;\xleftarrow{\partial J/\partial \yhat}\;
  \yhat \;\xleftarrow{\partial \yhat/\partial z^{(2)}}\;
  z^{(2)} \;\xleftarrow{\partial z^{(2)}/\partial h_i}\;
  h_i \;\xleftarrow{\partial h_i/\partial z^{(1)}_i}\;
  z^{(1)}_i & \;\xleftarrow{\partial z^{(1)}_i/\partial w^{(1)}_{ij}}\;
  w^{(1)}_{ij}                                                          \\[0.4em]
  J \;\xleftarrow{\partial J/\partial \yhat}\;
  \yhat \;\xleftarrow{\partial \yhat/\partial z^{(2)}}\;
  z^{(2)} \;\xleftarrow{\partial z^{(2)}/\partial h_i}\;
  h_i \;\xleftarrow{\partial h_i/\partial z^{(1)}_i}\;
  z^{(1)}_i & \;\xleftarrow{\partial z^{(1)}_i/\partial b^{(1)}_i}\;
  b^{(1)}_i
\end{align*}

As in the output layer, the weight and bias chains share every link except
the last.

We already know $\delta^{(2)} = \partial J / \partial z^{(2)} = -0.168$.
The chain rule gives the error signal at hidden unit $i$:

% ──────────────────────────────────────────────────────────────────────────────
\begin{patternbox}
  \textbf{Error signal at hidden unit $i$:}
  \[
    \delta^{(1)}_i
    \;=\; \pd{J}{z^{(1)}_i}
    \;=\; \delta^{(2)}
    \cdot w^{(2)}_i
    \cdot \sigmoid'\!\left(z^{(1)}_i\right).
  \]

  Three factors multiplied together:
  \begin{enumerate}[nosep]
    \item $\delta^{(2)}$: error signal from the layer ahead.
    \item $w^{(2)}_i$: how strongly $h_i$ connects to the output.
    \item $\sigmoid'(z^{(1)}_i)$: how sensitive $h_i$ is to its own
          pre-activation.
  \end{enumerate}

  Once $\delta^{(1)}_i$ is known, the weight gradients follow the same rule
  as before:
  \[
    \pd{J}{w^{(1)}_{ij}} = \delta^{(1)}_i \cdot x_j, \qquad
    \pd{J}{b^{(1)}_i}    = \delta^{(1)}_i.
  \]
\end{patternbox}

\vspace{0.8em}

\begin{figure}[H]
  \centering
  \begin{tikzpicture}[
      preact/.style={circle, draw=mainblue, fill=white,
          minimum size=1.1cm, thick, font=\small},
      param/.style={rectangle, rounded corners=3pt, draw=orange!80!black,
          fill=notebg, thick, minimum size=0.75cm, font=\small},
      linklabel/.style={font=\small, orange!60!black},
      >=Stealth
    ]
    % ════ Hidden unit 1 ═══════════════════════════════════════════════════════
    \begin{scope}
      \node[preact] (z1) at (6, 0) {$z^{(1)}_1$};
      \draw[->, thick] (1.8,  0.8) node[left, font=\small] {$x_1 = 0$} -- (z1);
      \draw[->, thick] (1.8, -0.8) node[left, font=\small] {$x_2 = 1$} -- (z1);

      \node[param] (w11) at (2.4,  2.5) {$w^{(1)}_{11}$};
      \node[param] (w12) at (2.4, -2.5) {$w^{(1)}_{12}$};
      \node[param] (b11) at (6.0,  2.5) {$b^{(1)}_1$};
      \draw[dotted, very thick, orange!80!black] (w11) --
      node[linklabel, above right, pos=0.4] {$\pd{z^{(1)}_1}{w^{(1)}_{11}}$} (z1);
      \draw[dotted, very thick, orange!80!black] (w12) --
      node[linklabel, below right, pos=0.4] {$\pd{z^{(1)}_1}{w^{(1)}_{12}}$} (z1);
      \draw[dotted, very thick, orange!80!black] (b11) --
      node[linklabel, right, pos=0.45, xshift=2pt] {$\pd{z^{(1)}_1}{b^{(1)}_1}$} (z1);

      \draw[->, thick, red!70!black, dashed] (11.4, 0) --
      node[above, font=\small] {$\delta^{(1)}_1 = \pd{J}{z^{(1)}_1}$}
      node[below, font=\scriptsize, gray!80!black, yshift=-2pt]
        {from Figure~\ref{fig:hidden-chain}} (z1);
    \end{scope}

    % ── Separator ──────────────────────────────────────────────────────────────
    \draw[gray!50, dashed] (0, -3.3) -- (11.6, -3.3);

    % ════ Hidden unit 2 ═══════════════════════════════════════════════════════
    \begin{scope}[yshift=-6.6cm]
      \node[preact] (z2) at (6, 0) {$z^{(1)}_2$};
      \draw[->, thick] (1.8,  0.8) node[left, font=\small] {$x_1 = 0$} -- (z2);
      \draw[->, thick] (1.8, -0.8) node[left, font=\small] {$x_2 = 1$} -- (z2);

      \node[param] (w21) at (2.4,  2.5) {$w^{(1)}_{21}$};
      \node[param] (w22) at (2.4, -2.5) {$w^{(1)}_{22}$};
      \node[param] (b12) at (6.0,  2.5) {$b^{(1)}_2$};
      \draw[dotted, very thick, orange!80!black] (w21) --
      node[linklabel, above right, pos=0.4] {$\pd{z^{(1)}_2}{w^{(1)}_{21}}$} (z2);
      \draw[dotted, very thick, orange!80!black] (w22) --
      node[linklabel, below right, pos=0.4] {$\pd{z^{(1)}_2}{w^{(1)}_{22}}$} (z2);
      \draw[dotted, very thick, orange!80!black] (b12) --
      node[linklabel, right, pos=0.45, xshift=2pt] {$\pd{z^{(1)}_2}{b^{(1)}_2}$} (z2);

      \draw[->, thick, red!70!black, dashed] (11.4, 0) --
      node[above, font=\small] {$\delta^{(1)}_2 = \pd{J}{z^{(1)}_2}$}
      node[below, font=\scriptsize, gray!80!black, yshift=-2pt]
        {from Figure~\ref{fig:hidden-chain}} (z2);
    \end{scope}
  \end{tikzpicture}
  \caption{Zooming in on the hidden pre-activations.  The hidden layer has
    four weights and two biases, each linked to one $z^{(1)}_i$.  As in the
    output layer, each parameter's gradient is $\delta^{(1)}_i$ times the
    label on its link.}
  \label{fig:hidden-local}
\end{figure}

% ──────────────────────────────────────────────────────────────────────────────
\subsection{Applying the Pattern to $h_1$}

\[
  \delta^{(1)}_1
  = \delta^{(2)} \cdot w^{(2)}_1 \cdot \sigmoid'(z^{(1)}_1)
  = (-0.168)(0.5)(0.250)
  = \sol[3cm]{-0.021}
\]

Because $z^{(1)}_1 = w^{(1)}_{11} x_1 + w^{(1)}_{12} x_2 + b^{(1)}_1$,
the local derivatives are:
\[
  \pd{z^{(1)}_1}{w^{(1)}_{11}} = x_1 = 0, \qquad
  \pd{z^{(1)}_1}{w^{(1)}_{12}} = x_2 = 1, \qquad
  \pd{z^{(1)}_1}{b^{(1)}_1} = 1.
\]

The gradients for row 1 of $\bW^{(1)}$ and for $b^{(1)}_1$:

\begin{align*}
  \pd{J}{w^{(1)}_{11}} & = \delta^{(1)}_1 \cdot x_1 = (-0.021)(0) = \sol[2.5cm]{0.000}  \\[4pt]
  \pd{J}{w^{(1)}_{12}} & = \delta^{(1)}_1 \cdot x_2 = (-0.021)(1) = \sol[2.5cm]{-0.021} \\[4pt]
  \pd{J}{b^{(1)}_1}    & = \delta^{(1)}_1 \cdot 1 = \sol[2.5cm]{-0.021}
\end{align*}

\begin{notebox}
  $\partial J / \partial w^{(1)}_{11} = 0$ because $x_1 = 0$: this input was
  zero, so the weight connecting it had no effect on the prediction --- and
  therefore no effect on the loss.
\end{notebox}

% ──────────────────────────────────────────────────────────────────────────────
\subsection{Applying the Pattern to $h_2$}

\[
  \delta^{(1)}_2
  = \delta^{(2)} \cdot w^{(2)}_2 \cdot \sigmoid'(z^{(1)}_2)
  = (-0.168)(0.5)(0.235)
  \approx \sol[3cm]{-0.020}
\]

Because $z^{(1)}_2 = w^{(1)}_{21} x_1 + w^{(1)}_{22} x_2 + b^{(1)}_2$,
the local derivatives are:
\[
  \pd{z^{(1)}_2}{w^{(1)}_{21}} = x_1 = 0, \qquad
  \pd{z^{(1)}_2}{w^{(1)}_{22}} = x_2 = 1, \qquad
  \pd{z^{(1)}_2}{b^{(1)}_2} = 1.
\]

The gradients for row 2 of $\bW^{(1)}$ and for $b^{(1)}_2$:

\begin{align*}
  \pd{J}{w^{(1)}_{21}} & = \delta^{(1)}_2 \cdot x_1 = (-0.020)(0) = \sol[2.5cm]{0.000}  \\[4pt]
  \pd{J}{w^{(1)}_{22}} & = \delta^{(1)}_2 \cdot x_2 = (-0.020)(1) = \sol[2.5cm]{-0.020} \\[4pt]
  \pd{J}{b^{(1)}_2}    & = \delta^{(1)}_2 \cdot 1 = \sol[2.5cm]{-0.020}
\end{align*}

% ══════════════════════════════════════════════════════════════════════════════
\section{Summary of All Gradients}

\begin{center}
  \renewcommand{\arraystretch}{1.5}
  \begin{tabular}{lrr}
    \toprule
    Parameter      & Current value    & Gradient $\partial J / \partial(\cdot)$ \\
    \midrule
    $w^{(2)}_1$    & $\phantom{-}0.5$ & $-0.084$                                \\
    $w^{(2)}_2$    & $\phantom{-}0.5$ & $-0.105$                                \\
    $b^{(2)}$      & $\phantom{-}0.0$ & $-0.168$                                \\
    \midrule
    $w^{(1)}_{11}$ & $\phantom{-}0.5$ & $\phantom{-}0.000$                      \\
    $w^{(1)}_{12}$ & $-0.5$           & $-0.021$                                \\
    $b^{(1)}_1$    & $\phantom{-}0.5$ & $-0.021$                                \\
    \midrule
    $w^{(1)}_{21}$ & $\phantom{-}0.5$ & $\phantom{-}0.000$                      \\
    $w^{(1)}_{22}$ & $\phantom{-}0.5$ & $-0.020$                                \\
    $b^{(1)}_2$    & $\phantom{-}0.0$ & $-0.020$                                \\
    \bottomrule
  \end{tabular}
\end{center}
All non-zero gradients are \textbf{negative}, so gradient descent is expected to move those
parameters up.
The two weights attached to $x_1 = 0$ ($w^{(1)}_{11}$ and $w^{(1)}_{21}$)
have zero gradient and should stay put.
% ══════════════════════════════════════════════════════════════════════════════
\section{Update}

With all gradients in hand, we apply one gradient descent step
(Section 3.2) to every parameter:
\[
  w \;\longleftarrow\; w - \eta \,\pd{J}{w},
  \qquad \eta = 0.1.
\]

\begin{center}
  \renewcommand{\arraystretch}{1.5}
  \begin{tabular}{lrlr}
    \toprule
    Parameter      & Old              & $\text{old} - \eta \cdot \text{gradient}$ & New                   \\
    \midrule
    $w^{(2)}_1$    & $\phantom{-}0.5$ & $0.5 - 0.1(-0.084)$                       & \sol[1.6cm]{$0.508$}  \\
    $w^{(2)}_2$    & $\phantom{-}0.5$ & $0.5 - 0.1(-0.105)$                       & \sol[1.6cm]{$0.510$}  \\
    $b^{(2)}$      & $\phantom{-}0.0$ & $0.0 - 0.1(-0.168)$                       & \sol[1.6cm]{$0.017$}  \\
    \midrule
    $w^{(1)}_{11}$ & $\phantom{-}0.5$ & $0.5 - 0.1(0.000)$                        & \sol[1.6cm]{$0.500$}  \\
    $w^{(1)}_{12}$ & $-0.5$           & $-0.5 - 0.1(-0.021)$                      & \sol[1.6cm]{$-0.498$} \\
    $b^{(1)}_1$    & $\phantom{-}0.5$ & $0.5 - 0.1(-0.021)$                       & \sol[1.6cm]{$0.502$}  \\
    \midrule
    $w^{(1)}_{21}$ & $\phantom{-}0.5$ & $0.5 - 0.1(0.000)$                        & \sol[1.6cm]{$0.500$}  \\
    $w^{(1)}_{22}$ & $\phantom{-}0.5$ & $0.5 - 0.1(-0.020)$                       & \sol[1.6cm]{$0.502$}  \\
    $b^{(1)}_2$    & $\phantom{-}0.0$ & $0.0 - 0.1(-0.020)$                       & \sol[1.6cm]{$0.002$}  \\
    \bottomrule
  \end{tabular}
\end{center}

Running the forward pass again on the same input $(x_1, x_2) = (0, 1)$
with the new parameters:

\begin{center}
  \renewcommand{\arraystretch}{1.4}
  \begin{tabular}{lcc}
    \toprule
    Quantity           & Before update & After update         \\
    \midrule
    Prediction $\yhat$ & $0.637$       & \sol[1.6cm]{$0.643$} \\
    Loss $J$           & $0.132$       & \sol[1.6cm]{$0.127$} \\
    \bottomrule
  \end{tabular}
\end{center}


After one update step the prediction has moved closer to the target $y = 1$,
and the loss has dropped.

\begin{warnbox}
  \textbf{One forward--backward pass updates weights for one input.}
  In practice we cycle through all four XOR inputs (or use mini-batches),
  accumulate gradients, and take one update step per batch.
  After hundreds or thousands of such cycles the network converges and
  correctly classifies all four XOR inputs.
\end{warnbox}

% ══════════════════════════════════════════════════════════════════════════════
\section{The Big Picture}

The full training loop is now complete in principle:

\begin{center}
  \begin{tikzpicture}[
      box/.style={draw=mainblue, fill=boxbg, rounded corners=4pt,
          minimum width=3.8cm, minimum height=1.4cm,
          font=\normalsize, align=center, thick},
      flow/.style={->, thick, mainblue},
      >=Stealth
    ]
    \node[box] (fwd)  at (0,    0)   {Forward pass\\(Section 3.3)};
    \node[box] (loss) at (5.2,  0)   {Compute loss\\$J = (y-\yhat)^2$};
    \node[box] (bwd)  at (10.4, 0)   {Backprop\\(Section 3.4)};
    \node[box] (upd)  at (5.2, -3.0) {Gradient descent\\update};

    \draw[flow] (fwd)  -- node[above, font=\small, black] {$\yhat$} (loss);
    \draw[flow] (loss) -- node[above, font=\small, black] {$J$} (bwd);
    \draw[flow] (bwd.south) |- node[pos=0.25, right, font=\small, black]
      {$\pd{J}{w},\ \pd{J}{b}$} (upd.east);
    \draw[flow] (upd.west) -| node[pos=0.75, left, font=\small, black]
      {new $w,\, b$} (fwd.south);
  \end{tikzpicture}
\end{center}

\vspace{0.5em}

\begin{notebox}
  \textbf{Note:} Plain gradient descent works, but it can be slow or sensitive to the
  learning rate. ``Adam" is a smarter update rule that adapts the step size for
  each weight automatically.
\end{notebox}

\end{document}
